\subsubsection{Attention while driving}

For the driving simulator experiment Gaspard Lefort and I helped set up the
experiment, it was needed to install the software for recording the screen,
setup an eye tracker and a EEG cap. We had help from the mechatronics
departement to have a EEG cap of the Emotiv brand, they could make us try it
with a software where we could command the movement of an object with our mind.
When the experiment was set up, we were here for the experiment on the first
subjects to make sure everything worked and have a protocol on how to make the
experiment and compile the results.

The experiment asked the subject to have an eye tracker and an EEG cap while
driving on a driving simulator, during the driving some events were triggered by
the researcher, rainfall, phone notification, night, to disturb the subject
while driving. The idea was to see with the eye tracker if the subject looked at
the phone when the notification occured and the EEG cap would record the shift
of attention. The driving simulator software gave us all the details of the run,
when the events were triggered, when the subject went over the speed limit,
etc... The setup is shown in \cref{fig:driving-simulator}.

\insertfigure
    {assets/images/driving simulator-1.jpg}
    {width=0.7\textwidth}
    {Driving Simulator Setup (without the EEG)}
    {fig:driving-simulator}

We were not here to conduct most of the experiments but I helped B.Tech interns
to compile the data, the idea was to have a timeline where we could plot all the
events, the processed attention based on the EEG, and the gaze (if it was
centered or not). I gave idea on how to process the gaze with a computer vision
algorithm because it was only recorded as a hotspot on the video capture of the
screen during the run.

The experiments were not done by the end of the internship, results will be
missing here.